% !TEX root = ../../main.tex

\section{Integrality}\label{sec:integrality}

Integrality of the limit varifold is an essential input to the regularity argument: it ensures that tangent cones at singular points are integral cones, which constrains them to the junction configurations (wedges, folds, and higher-order junctions) needed to derive a contradiction when the $\alpha$-structural hypothesis (Definition~\ref{defn:alpha-structural}) fails (see Section~\ref{subsec:alpha-structural}). We begin by recalling the following criterion from de~Lellis--Tasnady~\cite{DLT13}, used in the no-mass-cancellation case (Section~\ref{subsec:integrality-no-cancel}):

\begin{proposition}[Varifolds vs.\ Caccioppoli set limits~{\cite[Proposition~A.1]{DLT13}}]\label{prop:p2-DLT-integrality}
    Let $\{\Omega^k\}$ be a sequence of Caccioppoli sets and $U$ an open subset of $M$. Assume that
    \begin{enumerate}[label=\textup{(\roman*)}]
        \item $D\chi_{\Omega^k} \to D\chi_\Omega$ in the sense of measures in $U$;
        \item $\mathrm{Per}(\Omega^k, U) \to \mathrm{Per}(\Omega, U)$
    \end{enumerate}
    for some Caccioppoli set $\Omega$. Then the varifolds $|\partial^*\Omega^k|$ converge to $|\partial^*\Omega|$ in the sense of varifolds.
\end{proposition}

The main result of this section is:

\begin{theorem}[Integrality of the limit varifold]\label{thm:integrality}
Let $(M^{n+1}, g)$ be a closed Riemannian manifold with $n \geq 2$, and let $\Phi$ be a non-excessive ONVP sweepout with width $W$. Let $x_i \to x_0 \in \mathfrak{m}(\Phi)$ be a min-max sequence and $V$ the varifold limit of $|\partial^*\Omega(x_i)|$.
\begin{enumerate}[label=\textup{(\alph*)}]
    \item If no mass cancellation occurs at $x_0$ (i.e., $\mathbf{M}(\Phi(x_0)) = W$), then $V = |\partial^*\Omega(x_0)|$ is an integral $n$-varifold.
    \item If mass cancellation occurs at $x_0$, then $\|V\|(A) \geq \mathcal{H}^n(\partial^*\Omega(x_0) \cap A)$ for every Borel $A$; moreover the restriction of $V$ to any open neighborhood of $\partial^*\Omega(x_0)$ that is disjoint from $\spt(\|V\| - |D\chi_{\Omega(x_0)}|)$ is integral. Integrality of $V$ on the full support $\spt\|V\|$ is conditional on Proposition~\ref{prop:positive-density} (positive density of the varifold limit $\|V\|$-a.e.\ on the cancelled mass $\mu_{\mathrm{c}} = \|V\| - |D\chi_{\Omega(x_0)}|$); see Remark~\ref{rem:cancelled-mass-gap}.
\end{enumerate}
\end{theorem}

Part~(a) is the de~Lellis--Tasnady criterion (Theorem~\ref{thm:integrality-mult-one} below); part~(b) is Lemma~\ref{lem:density-lower-bound} combined with the sheet decomposition (Lemma~\ref{lem:sheet-decomposition}) applied to $\mathcal{H}^n$-a.e.\ reduced boundary point, and with the gap in the cancelled mass spelled out in Remark~\ref{rem:cancelled-mass-gap}. Under the multiplicity-one conjecture (no mass cancellation), part~(a) applies and integrality is unconditional.

\subsection{The no-mass-cancellation case}\label{subsec:integrality-no-cancel}

\begin{theorem}[Integrality without mass cancellation]\label{thm:integrality-mult-one}
Assume no mass cancellation occurs at $x_0$, i.e., $\mathbf{M}(\Phi(x_0)) = W$. Then $V = |\partial^*\Omega(x_0)|$; in particular, $V$ is an integral varifold.
\end{theorem}

\begin{proof}
Since $\Phi$ is continuous in the flat topology, we have $\Omega(x_i) \to \Omega(x_0)$ in $L^1(M)$, so $D\chi_{\Omega(x_i)} \to D\chi_{\Omega(x_0)}$ weakly as measures. Moreover, the no-mass-cancellation hypothesis gives
\[
\mathrm{Per}(\Omega(x_i)) = \mathbf{M}(\Phi(x_i)) \to W = \mathbf{M}(\Phi(x_0)) = \mathrm{Per}(\Omega(x_0)).
\]
Both conditions of Proposition~\ref{prop:p2-DLT-integrality} are satisfied, so $V = |\partial^*\Omega(x_0)|$ is the integral varifold induced by the Caccioppoli boundary $\partial^*\Omega(x_0)$.
\end{proof}

\begin{figure}[ht]
    \centering
    \begin{tikzpicture}[scale=1.0]
      \def\R{2.2}
      \def\h{0.9}
      \pgfmathsetmacro{\a}{sqrt(\R*\R - \h*\h)}
      \def\ew{0.4}
      \pgfmathsetmacro{\angR}{atan2(\h, \a)}

      % Fill cap (Omega): front ellipse arc + sphere arc over top
      \fill[fillA, opacity=0.4]
        (-\a, \h) arc (180:360:{\a} and {\ew})
        arc (\angR:180-\angR:\R);

      % Sphere outline
      \draw[thick] (0,0) circle (\R);

      % Boundary curve — front (solid)
      \draw[very thick, bindA] (-\a, \h) arc (180:360:{\a} and {\ew});
      % Boundary curve — back (dashed)
      \draw[thick, bindA, dashed] (-\a, \h) arc (180:0:{\a} and {\ew});

      % Equator (V) — front (solid)
      \draw[very thick, bindB] (-\R, 0) arc (180:360:{\R} and {0.45});
      % Equator (V) — back (dashed)
      \draw[thick, bindB, dashed] (-\R, 0) arc (180:0:{\R} and {0.45});

      % Labels
      \node[bindA] at (0, 1.75) {$\Omega(x_i)$};
      \node[bindA, right] at (\a+0.1, \h-0.15) {$|\partial^*\!\Omega(x_i)|$};
      \node[bindB, right] at (\R+0.1, -0.15) {$V = |\partial^*\!\Omega(x_0)|$};
      \node[above left] at ({-\R*cos(40)}, {\R*sin(40)}) {$M$};
    \end{tikzpicture}
    \caption{The no-mass-cancellation case on a closed manifold $M$. The shaded region is the Caccioppoli set $\Omega(x_i)$, with reduced boundary $|\partial^*\!\Omega(x_i)|$ (red). As $x_i \to x_0$, the varifolds $|\partial^*\!\Omega(x_i)|$ converge to $V = |\partial^*\!\Omega(x_0)|$ (blue), which is integral with multiplicity~$1$.}
    \label{fig:caccioppoli-set}
\end{figure}

The identification $V = |\partial^*\Omega(x_0)|$ has two important consequences beyond integrality: it tells us that $V$ is the boundary of a Caccioppoli set, which constrains the possible tangent cones and is used to verify the $\alpha$-structural hypothesis (Theorem~\ref{thm:alpha-hypo}); and it connects $V$ to the non-excessive structure of the sweepout, which provides the homotopic minimizing property and hence stability (Proposition~\ref{prop:stability}). Together with integrality, these are the inputs to Wickramasekera's regularity theorem (Section~\ref{sec:regularity}).

\subsection{The mass cancellation case}\label{subsec:integrality-cancel}

When mass cancellation occurs ($\mathbf{M}(\Phi(x_0)) < W$), we still have $\Omega(x_i) \to \Omega(x_0)$ in $L^1$, but the perimeter convergence condition in Proposition~\ref{prop:p2-DLT-integrality} fails:
\[
    \mathrm{Per}(\Omega(x_i)) \to W > \mathrm{Per}(\Omega(x_0)).
\]
The mass defect $W - \mathrm{Per}(\Omega(x_0)) > 0$ is carried by sheets of $\partial^*\Omega(x_i)$ that converge in pairs, cancelling in the flat topology but contributing positively to varifold mass. In particular, $V \neq |\partial^*\Omega(x_0)|$, and the integrality of $V$ does not follow from Proposition~\ref{prop:p2-DLT-integrality}.

The analysis splits $\|V\|$ into two parts via the comparison
\begin{equation}\label{eq:V-vs-D-chi-decomp}
    \|V\| = |D\chi_{\Omega(x_0)}| + \mu_{\mathrm{c}}, \qquad \mu_{\mathrm{c}} := \|V\| - |D\chi_{\Omega(x_0)}| \geq 0,
\end{equation}
where the non-negativity $\mu_{\mathrm{c}} \geq 0$ is established in Lemma~\ref{lem:density-lower-bound} by lower semicontinuity of total variation. The first part, $|D\chi_{\Omega(x_0)}|$, is the reduced-boundary component: at each $p \in \partial^*\Omega(x_0)$ density is $\geq 1$, and the sheet decomposition (Lemma~\ref{lem:sheet-decomposition}) further refines this to integer multiplicity. The second part, $\mu_{\mathrm{c}}$ (the \emph{cancelled mass}), satisfies $\mu_{\mathrm{c}}(M) = W - \mathrm{Per}(\Omega(x_0)) > 0$, is supported on the cancellation set where pairs of sheets of $\partial^*\Omega(x_i)$ converge to a common location; establishing its rectifiability and integer multiplicity is a separate problem discussed in Remark~\ref{rem:cancelled-mass-gap}.

The main tool for the $\partial^*\Omega(x_0)$-part is the following graphical decomposition lemma, which we state first.


\begin{lemma}[Lower bound on the limit mass by the flat limit]\label{lem:density-lower-bound}
Let $(M^{n+1}, g)$ be a closed Riemannian manifold with $n \geq 2$, and let $\Phi$ be a (not necessarily non-excessive) ONVP sweepout. Let $V$ be the varifold limit of $|\partial^*\Omega(x_i)|$ along a min-max sequence $x_i \to x_0 \in \mathfrak{m}(\Phi)$. Then
\begin{equation}\label{eq:V-dominates-DChi}
    \|V\|(A) \geq |D\chi_{\Omega(x_0)}|(A) = \mathcal{H}^n(\partial^*\Omega(x_0) \cap A)
    \quad \text{for every Borel } A \subset M.
\end{equation}
In particular, $\Theta(\|V\|, p) \geq 1$ for every $p \in \partial^*\Omega(x_0)$.
\end{lemma}

The inequality~\eqref{eq:V-dominates-DChi} follows from the lower semicontinuity of total variation under $L^1$ convergence, combined with the varifold convergence $|\partial^*\Omega(x_i)| \to V$; non-excessiveness is not used. The lemma gives positive density on $\partial^*\Omega(x_0) \cap \spt\|V\|$, on which the sheet decomposition (Lemma~\ref{lem:sheet-decomposition}) yields integer multiplicity. The complementary cancelled mass $\mu_{\mathrm{c}}$ from~\eqref{eq:V-vs-D-chi-decomp} is treated in Remark~\ref{rem:cancelled-mass-gap}.

% [Commented out: old proof by contradiction via isoperimetric replacement]
% \begin{proof}
% Since $V$ is stationary, the density $\Theta(\|V\|, p)$ exists at every point (Proposition~\ref{prop:p2-monotonicity}). We argue by contradiction: suppose there exists $p \in \spt\|V\|$ with $\Theta(\|V\|, p) = 0$. The strategy is as follows: the vanishing density forces $\Omega(x_i)$ to be nearly empty or nearly full inside small balls around~$p$. An isoperimetric replacement then produces a Caccioppoli set $E_i$ with strictly less perimeter, from which we construct an $I$-replacement family contradicting the non-excessiveness of $\Phi$.

% \medskip\noindent\textbf{Step~1} (Vanishing mass at small scales).
% Since the modified density ratio $e^{Cr}\,\omega_n^{-1} r^{-n}\|V\|(B_r(p))$ is non-decreasing in $r$ (Proposition~\ref{prop:p2-monotonicity}) and its limit as $r \to 0$ equals $\Theta(\|V\|, p) = 0$, we have
% \begin{gather}
%     \|V\|(B_r(p)) \leq \omega_n\, \alpha(r)\, r^n, \label{eq:density-vanishing} \\
%     \alpha(r) := e^{Cr}\,\omega_n^{-1} r^{-n}\|V\|(B_r(p)) \xrightarrow{r \to 0} 0. \notag
% \end{gather}
% In particular, $\|V\|(B_r(p)) = o(r^n)$. Now fix $\varepsilon > 0$ and choose $r_0 > 0$ so that $\alpha(r) < \varepsilon$ for all $r < r_0$. Since $r \mapsto \|V\|(B_r(p))$ is non-decreasing and hence has at most countably many jumps, $\|V\|(\partial B_r(p)) = 0$ for all but countably many $r$. For any such $r < r_0$, the Portmanteau theorem (as in~\eqref{eq:tilt-excess-transfer}) applied to the indicator $\mathbf{1}_{B_r(p)}$ gives
% \begin{equation}\label{eq:mass-transfer}
%     \mathrm{Per}(\Omega(x_i), B_r(p)) = \|\partial^*\Omega(x_i)\|(B_r(p)) \to \|V\|(B_r(p)) \quad \text{as } i \to \infty.
% \end{equation}
% Combining: for every $\varepsilon > 0$, there exists $r_0 > 0$ such that for each (good) $r < r_0$, there exists $i_0 = i_0(r, \varepsilon)$ with
% \begin{equation}\label{eq:per-small}
%     \mathrm{Per}(\Omega(x_i), B_r(p)) < \varepsilon\, \omega_n\, r^n \quad \text{for all } i \geq i_0.
% \end{equation}
%
% \medskip\noindent\textbf{Step~2} (Volume dichotomy via the isoperimetric inequality).
% Applying the relative isoperimetric inequality (Proposition~\ref{prop:relative-isoperimetric}) to $\Omega = \Omega(x_i)$ and using~\eqref{eq:per-small}:
% \begin{multline}\label{eq:volume-dichotomy}
%     \min\bigl(\mathcal{L}^{n+1}(\Omega(x_i) \cap B_r(p)),\;
%     \mathcal{L}^{n+1}(B_r(p) \setminus \Omega(x_i))\bigr) \\
%     \leq \bigl(C_I\,\varepsilon\,\omega_n\, r^n\bigr)^{(n+1)/n}
%     = C\,\varepsilon^{(n+1)/n}\, r^{n+1}
% \end{multline}
% for each good $r < \min(r_0, r_1)$ and all $i \geq i_0(r, \varepsilon)$. Since $\varepsilon > 0$ is arbitrary, the smaller side has volume $o(r^{n+1})$ as $r \to 0$ then $i \to \infty$: the set $\Omega(x_i) \cap B_r(p)$ is either nearly empty or nearly full.
%
% \medskip\noindent\textbf{Step~3} (Replacement and contradiction with non-excessiveness; cf.~\cite[Proposition~5.1]{DLT13}).
% By passing to a subsequence, we may assume $x_i \nearrow x_0$, so that $x_0 \in \mathfrak{m}_L(\Phi)$ and $x_0$ is not left excessive (Theorem~\ref{thm:non-excessive-existence}). By Step~2, passing to a further subsequence, either
% $\mathcal{L}^{n+1}(\Omega(x_i) \cap B_r(p)) = o(r^{n+1})$ (nearly empty) or
% $\mathcal{L}^{n+1}(B_r(p) \setminus \Omega(x_i)) = o(r^{n+1})$ (nearly full).
% Let $f_i(\rho)$ denote the smaller volume:
% \begin{equation}\label{eq:volume-small}
%     f_i(r) := \min\bigl(\mathcal{L}^{n+1}(\Omega(x_i) \cap B_r(p)),\;\mathcal{L}^{n+1}(B_r(p) \setminus \Omega(x_i))\bigr) = o(r^{n+1}).
% \end{equation}
%
% \emph{Finding a good radius.}
% The goal is to find a radius $\rho_i$ at which the boundary $\partial B_{\rho_i}(p)$ slices off less of $S_i$ than the isoperimetric inequality accounts for, so that we can construct a cheaper replacement.
%
% Let $S_i$ denote the side achieving the minimum in~\eqref{eq:volume-small}: $S_i = \Omega(x_i)$ if $\mathcal{L}^{n+1}(\Omega(x_i) \cap B_r(p)) \leq \mathcal{L}^{n+1}(B_r(p) \setminus \Omega(x_i))$, and $S_i = B_r(p) \setminus \Omega(x_i)$ otherwise, so that $f_i(\rho) = \mathcal{L}^{n+1}(S_i \cap B_\rho(p))$. The coarea formula gives
% \[
%     f_i'(\rho) = \mathcal{H}^n(\partial B_\rho(p) \cap S_i) \quad \text{for a.e.\ } \rho \in (0, r).
% \]
% Meanwhile, the relative isoperimetric inequality (Proposition~\ref{prop:relative-isoperimetric}) gives
% \begin{equation}\label{eq:isoperimetric-lower}
%     \mathrm{Per}(\Omega(x_i), B_\rho(p)) \geq c_n\, f_i(\rho)^{n/(n+1)},
% \end{equation}
% where $c_n = c_n(M, g, n) > 0$. We claim that there exists $\rho_i \in (r/2, r)$ with $f_i(\rho_i) > 0$ and
% \begin{equation}\label{eq:good-radius}
%     f_i'(\rho_i) < c_n\, f_i(\rho_i)^{n/(n+1)} \leq \mathrm{Per}(\Omega(x_i), B_{\rho_i}(p)).
% \end{equation}
% The first inequality in~\eqref{eq:good-radius} says that the ``slice cost'' $\mathcal{H}^n(\partial B_{\rho_i} \cap S_i)$ is strictly less than the perimeter inside $B_{\rho_i}$; at such a radius, replacing $\Omega(x_i)$ inside $B_{\rho_i}$ by filling in (or emptying out) $S_i$ reduces the total perimeter.
%
% To prove the claim, note that $p \in \spt\|V\|$ gives $\|V\|(B_{r/2}(p)) > 0$, so by~\eqref{eq:mass-transfer}, $\mathrm{Per}(\Omega(x_i), B_{r/2}(p)) > 0$ for $i$ large. This forces $f_i(\rho) > 0$ for all $\rho \geq r/2$: indeed, $\mathrm{Per} > 0$ gives a point $q \in \partial^*\Omega(x_i) \cap B_{r/2}(p)$, and at any reduced boundary point the set $\Omega(x_i)$ has $\mathcal{L}^{n+1}$-density $1/2$, so both $\Omega(x_i)$ and its complement have positive volume in $B_\rho(p) \supset B_{r/2}(p) \ni q$. Now suppose for contradiction that $f_i'(\rho) \geq c_n\, f_i(\rho)^{n/(n+1)}$ for a.e.\ $\rho \in (r/2, r)$. By the ODE comparison argument of~\cite[Proposition~5.1, Step~2]{DLT13}, setting $\varphi = f_i^{1/(n+1)}$ and dividing:
% \[
%     \varphi'(\rho) = \frac{f_i'(\rho)}{(n{+}1)\,f_i(\rho)^{n/(n+1)}} \geq \frac{c_n}{n{+}1} \quad \text{a.e.\ on } (r/2, r).
% \]
% Integrating over $(r/2, r)$:
% \[
%     f_i(r)^{1/(n+1)} \geq f_i(r/2)^{1/(n+1)} + \frac{c_n\, r}{2(n{+}1)} \geq \frac{c_n\, r}{2(n{+}1)},
% \]
% hence $f_i(r) \geq (c_n / (2(n{+}1)))^{n+1}\, r^{n+1}$. Since $(c_n/(2(n{+}1)))^{n+1} > 0$ is independent of $r$ and $i$, this contradicts $f_i(r) = o(r^{n+1})$ from~\eqref{eq:volume-small}, establishing the claim.
%
% \emph{Construction of the replacement.}
% Define the Caccioppoli set
% \[
%     E_i := \Omega(x_i) \triangle (S_i \cap \overline{B}_{\rho_i}(p)),
% \]
% which removes the small side $S_i$ from $\Omega(x_i)$ inside $B_{\rho_i}(p)$: when $S_i = \Omega(x_i)$, this gives $E_i = \Omega(x_i) \setminus \overline{B}_{\rho_i}$; when $S_i = B_r \setminus \Omega(x_i)$, it gives $E_i = \Omega(x_i) \cup \overline{B}_{\rho_i}$. In both cases, $E_i$ agrees with $\Omega(x_i)$ outside $B_{\rho_i}(p)$, while inside $B_{\rho_i}(p)$ the boundary $\partial^*\Omega(x_i) \cap B_{\rho_i}$ is replaced by the spherical slice $\partial B_{\rho_i} \cap S_i$ (see Figure~\ref{fig:isoperimetric-replacement}). For a.e.\ $\rho_i$ (chosen so that $\|V\|(\partial B_{\rho_i}(p)) = 0$ and $\mathcal{H}^{n-1}(\partial^*\Omega(x_i) \cap \partial B_{\rho_i}(p)) = 0$):
% \begin{align}
%     \mathrm{Per}(E_i, B_r(p))
%     &= \mathrm{Per}(\Omega(x_i),\, B_r(p) \setminus \overline{B}_{\rho_i})
%        + f_i'(\rho_i), \label{eq:per-replacement} \\
%     \mathrm{Per}(\Omega(x_i), B_r(p))
%     &= \mathrm{Per}(\Omega(x_i),\, B_r(p) \setminus \overline{B}_{\rho_i})
%        + \mathrm{Per}(\Omega(x_i), B_{\rho_i}(p)). \label{eq:per-original}
% \end{align}
% Therefore:
% \begin{alignat}{2}
%     &\mathrm{Per}(E_i) - \mathrm{Per}(\Omega(x_i)) \notag \\
%     ={}&\mathrm{Per}(E_i, B_r(p)) - \mathrm{Per}(\Omega(x_i), B_r(p))
%        &\qquad &\text{($E_i = \Omega(x_i)$ outside $B_r$)} \notag \\
%     ={}&f_i'(\rho_i) - \mathrm{Per}(\Omega(x_i), B_{\rho_i}(p))
%        &\qquad &\text{(by~\eqref{eq:per-replacement} and~\eqref{eq:per-original})} \notag \\
%     <{}&0.
%        &\qquad &\text{(by~\eqref{eq:good-radius})} \label{eq:per-decrease}
% \end{alignat}
%
% \begin{figure}[ht]
%     ... (figure omitted in comments)
% \end{figure}
%
% [Commented out: Cases 1, 2a, 2b setup — superseded by direct 2b(i)/2b(ii) dichotomy]
% \emph{Contradiction with non-excessiveness.}
% We have $\mathrm{Per}(E_i) < \mathrm{Per}(\Omega(x_i)) \leq W$ from~\eqref{eq:per-decrease}, and $\mathrm{Per}(\Omega(x_0)) < W$ by the mass cancellation hypothesis. The goal is to construct a flat-continuous $I$-replacement on $I = (x_i, x_0]$ with $\limsup_{I \ni y \to x}\mathbf{M}(\Phi^I(y)) < W$ for every $x \in I$, making $x_0$ left excessive and contradicting non-excessiveness (Definition~\ref{def:excessive-interval}). Note that $x_i \notin I$, so no limsup condition is imposed at~$x_i$; the difficulty is that $\mathrm{Per}(\Omega(x_i))$ may equal~$W$. We distinguish two cases.
%
% \emph{Case~1: $\liminf_{i \to \infty}\mathrm{Per}(\Omega(x_i)) < W$.}
% Then there exist $\eta > 0$ and a subsequence (still denoted~$\{x_i\}$) with $\mathrm{Per}(\Omega(x_i)) \leq W - \eta$ for all~$i$.
% Since the sweepout is nested and $x_i < x_0$, we have $\Omega(x_i) \subset \Omega(x_0)$ with $\mathrm{Vol}(\Omega(x_0) \setminus \Omega(x_i)) \to 0$.
% Set $W_0 := \max\{W - \eta,\, \mathrm{Per}(\Omega(x_0))\} < W$ (independent of~$i$) and $\varepsilon := (W - W_0)/2 > 0$.
% Lemma~\ref{lem:interpolation} (applied with perimeter bound~$W$ and tolerance~$\varepsilon$) provides $\delta > 0$; for $i$ large enough, $\mathrm{Vol}(\Omega(x_0) \setminus \Omega(x_i)) \leq \delta$, so the lemma yields a nested $\mathcal{F}$-continuous family $\{\Omega_t\}_{t \in [0,1]}$ from $\Omega(x_i)$ to $\Omega(x_0)$ with
% \[
%     \mathrm{Per}(\Omega_t) \leq W_0 + \varepsilon = \frac{W + W_0}{2} < W
%     \quad\text{for all } t \in [0,1].
% \]
% Reparametrising this family on $\bar{I} = [x_i, x_0]$ gives an $I$-replacement with $\mathbf{M}(\Phi^I(x)) < W$ for all $x \in I = (x_i, x_0]$, making $x_0$ left excessive—a contradiction.
%
% \emph{Case~2: $\mathrm{Per}(\Omega(x_i)) \to W$.}
% The Interpolation Lemma alone cannot bring the perimeter below~$W$, since the uniform gap $W - \mathrm{Per}(\Omega(x_i)) \to 0$.
% We need an additional ingredient to decrease perimeter.
% \begin{list}{}{\leftmargin=1.5em \parsep=\parskip \topsep=\smallskipamount \itemsep=\medskipamount}
% \item \emph{Case~2a: $|\partial^*\Omega(x_i)|$ is stationary for infinitely many~$i$.}
% Passing to a subsequence we may assume $V_i := |\partial^*\Omega(x_i)|$ is stationary for every~$i$.
% Fix $R > 0$ smaller than the injectivity radius of $(M,g)$.
% Since $(M,g)$ is closed, the monotonicity formula for stationary varifolds on Riemannian manifolds (see e.g.~\cite[Section~17.6]{Sim83}) gives a constant $C = C(M,g) > 0$ such that for every stationary $n$-varifold $V$ on $(M,g)$, every $q \in \spt\|V\|$, and every $0 < \rho \leq R$,
% \[
%     \|V\|(B_\rho(q)) \geq e^{-C\rho}\,\omega_n\,\Theta^n(\|V\|, q)\,\rho^n.
% \]
% For each~$i$, the density of $V_i$ at every reduced boundary point $q \in \partial^*\Omega(x_i)$ satisfies $\Theta^n(\|V_i\|, q) = 1$ (since $V_i$ has multiplicity one). Setting $c := e^{-CR}\,\omega_n > 0$, which depends only on $(M,g)$, we obtain
% \begin{equation}\label{eq:mono-lower}
%     \|V_i\|(B_\rho(q)) \geq c\,\rho^n
%     \quad \text{for all } q \in \partial^*\Omega(x_i),\; 0 < \rho \leq R,\; i \in \mathbb{N}.
% \end{equation}
% Now pick $q_i \in \partial^*\Omega(x_i) \cap B_{r/2}(p)$ (which is non-empty for large $i$ since $p \in \spt\|V\|$). Applying~\eqref{eq:mono-lower} with $\rho = r/2$ gives $q_i \in B_{r/2}(p) \subset B_r(p)$, so
% \[
%     \|V_i\|(\overline{B}_{r}(p)) \geq \|V_i\|(B_{r/2}(q_i)) \geq c\,(r/2)^n.
% \]
% Since $V_i \to V$ as varifolds and $\overline{B}_{r}(p)$ is closed, the Portmanteau theorem yields
% \[
%     \|V\|(\overline{B}_{r}(p)) \geq \limsup_{i \to \infty} \|V_i\|(\overline{B}_{r}(p)) \geq c\,(r/2)^n.
% \]
% Dividing by $\omega_n r^n$ and letting $r \to 0$ gives $\Theta(\|V\|, p) \geq c/(2^n\omega_n) = e^{-CR}/2^n > 0$. No $I$-replacement is needed.
%
% \item \emph{Case~2b: $|\partial^*\Omega(x_i)|$ is not stationary for all large~$i$.}
% Passing to a subsequence, assume $|\partial^*\Omega(x_i)|$ is non-stationary for every~$i$.
% By the isoperimetric argument above, we have replacements $E_i$ with $\mathrm{Per}(E_i) < \mathrm{Per}(\Omega(x_i))$ and $E_i \mathbin{\triangle} \Omega(x_i) \subset \bar{B}_{\rho_i}(p)$. Passing to a subsequence, we may assume $E_i \subset \Omega(x_i)$ for all~$i$ (the nearly-empty case; the nearly-full case is symmetric). Then $\Omega(x_i)$ is not inner area-minimizing in $B_{\rho_i}(p)$.
% \begin{list}{}{\leftmargin=1.5em \parsep=\parskip \topsep=\smallskipamount \itemsep=\medskipamount}
% [End commented-out block]

\begin{proof}
We establish \eqref{eq:V-dominates-DChi} by a duality argument combining lower semicontinuity of total variation with varifold convergence, then deduce the density statement at reduced boundary points.

\medskip\noindent\textbf{Step 1} (The measure inequality \eqref{eq:V-dominates-DChi}).
Fix an open set $U \subset M$ and a function $\phi \in C_c(U)$ with $0 \leq \phi \leq 1$.
We first show
\begin{equation}\label{eq:phi-duality-comparison}
    \int_M \phi\, d|D\chi_{\Omega(x_0)}| \leq \int_M \phi\, d\|V\|.
\end{equation}

Let $\omega \in C^\infty_c(U; TM)$ with pointwise bound $|\omega(x)| \leq \phi(x)$ for every $x$.
By $\mathcal{F}$-continuity of $\Phi$ on $\bar I$, $\chi_{\Omega(x_i)} \to \chi_{\Omega(x_0)}$ in $L^1(M)$, hence
\[
    \int_M \chi_{\Omega(x_0)}\, \mathrm{div}\, \omega\, d\mathcal{L}^{n+1}
    = \lim_{i\to\infty} \int_M \chi_{\Omega(x_i)}\, \mathrm{div}\, \omega\, d\mathcal{L}^{n+1}.
\]
By definition of the distributional derivative, the left side equals $-\int_M \omega\cdot d(D\chi_{\Omega(x_0)})$ and each term on the right equals $-\int_M \omega\cdot d(D\chi_{\Omega(x_i)})$, so
\begin{equation}\label{eq:omega-limit}
    \int_M \omega\cdot d(D\chi_{\Omega(x_0)})
    = \lim_{i\to\infty} \int_M \omega\cdot d(D\chi_{\Omega(x_i)}).
\end{equation}
Writing $V_i = |\partial^*\Omega(x_i)|$ and noting $|D\chi_{\Omega(x_i)}|=\|V_i\|$ as Radon measures on $M$,
\[
    \left|\int_M \omega\cdot d(D\chi_{\Omega(x_i)})\right|
    \leq \int_M |\omega|\, d|D\chi_{\Omega(x_i)}|
    \leq \int_M \phi\, d\|V_i\|.
\]
Varifold convergence $V_i \to V$ on $\mathrm{Gr}_n(M)$ implies, in particular, $\|V_i\|\to\|V\|$ as Radon measures on $M$ tested against $\phi \in C_c(M)$ (take $f(x,S)=\phi(x)$ in Definition~\ref{def:p2-varifold-convergence}). Combining with \eqref{eq:omega-limit},
\begin{equation}\label{eq:sup-over-omega}
    \left|\int_M \omega\cdot d(D\chi_{\Omega(x_0)})\right|
    \leq \lim_{i\to\infty} \int_M \phi\, d\|V_i\|
    = \int_M \phi\, d\|V\|.
\end{equation}
By definition of the Radon measure $|D\chi_{\Omega(x_0)}|$ as the total variation of the vector-valued Radon measure $D\chi_{\Omega(x_0)}$,
\[
    \int_M \phi\, d|D\chi_{\Omega(x_0)}|
    = \sup_{\omega \in C^\infty_c(U;TM),\,|\omega|\leq \phi}
    \int_M \omega\cdot d(D\chi_{\Omega(x_0)}),
\]
so \eqref{eq:phi-duality-comparison} follows from \eqref{eq:sup-over-omega}.

Since \eqref{eq:phi-duality-comparison} holds for every $\phi \in C_c(U)$ with $0\leq\phi\leq 1$, outer regularity of the Radon measures $|D\chi_{\Omega(x_0)}|$ and $\|V\|$ yields $|D\chi_{\Omega(x_0)}|(U) \leq \|V\|(U)$, and inner regularity extends this to all Borel sets $A$, establishing~\eqref{eq:V-dominates-DChi}.

\medskip\noindent\textbf{Step 2} (Density at reduced boundary points).
By De~Giorgi's structure theorem, $|D\chi_{\Omega(x_0)}| = \mathcal{H}^n \llcorner \partial^*\Omega(x_0)$, and at every reduced boundary point $p$ the density equals~$1$ \cite[Corollary~15.8]{Mag12}:
\[
    \lim_{r\to 0} \frac{|D\chi_{\Omega(x_0)}|(B_r(p))}{\omega_n\, r^n} = 1.
\]
Combining with \eqref{eq:V-dominates-DChi} and the existence of $\Theta(\|V\|,p)$ via the monotonicity formula (Proposition~\ref{prop:p2-monotonicity}),
\[
    \Theta(\|V\|, p)
    = \lim_{r\to 0}\frac{\|V\|(B_r(p))}{\omega_n\, r^n}
    \geq \lim_{r\to 0}\frac{|D\chi_{\Omega(x_0)}|(B_r(p))}{\omega_n\, r^n}
    = 1
\]
for every $p \in \partial^*\Omega(x_0)$.
\end{proof}

\begin{remark}[The cancelled mass and an open gap]\label{rem:cancelled-mass-gap}
Lemma~\ref{lem:density-lower-bound} resolves the density lower bound on the reduced-boundary part $\partial^*\Omega(x_0) \cap \spt\|V\|$ but says nothing about the cancelled mass $\mu_{\mathrm{c}}$ from~\eqref{eq:V-vs-D-chi-decomp}. Three statements would close the cancellation case of integrality; we state them as questions and record what is and is not known.

\emph{(Q1) Is $\mu_{\mathrm{c}}$ mutually singular with $\mathcal{H}^n\llcorner\partial^*\Omega(x_0)$?}
The cancelled mass is carried by pairs of sheets of $\partial^*\Omega(x_i)$ approaching a common location from opposite sides, yielding no contribution to $D\chi_{\Omega(x_i)}$ in the flat limit. Heuristically these sheets accumulate on $\overline{\spt\|V\|} \setminus \partial^*\Omega(x_0)$, but a rigorous statement requires a quantitative compactness argument on the approximating boundaries. We do not prove this here.

\emph{(Q2) Is $\mu_{\mathrm{c}}$ $\mathcal{H}^n$-absolutely continuous with $\Theta^n(\mu_{\mathrm{c}}, p) \geq 1$ $\mu_{\mathrm{c}}$-a.e.?}
This is the analog of Lemma~\ref{lem:density-lower-bound} for the cancelled mass, and is the main open question in this paper. Resolving it affirmatively---equivalently, establishing Proposition~\ref{prop:positive-density}---closes the density lower bound $\Theta^n(\mu_{\mathrm{c}}, p) \geq c$ at $\|V\|$-a.e.\ point, from which rectifiability of $\mu_{\mathrm{c}}$ follows by Allard's theorem (Proposition~\ref{prop:allard-rectifiability}).

\emph{(Q3) If (Q1)--(Q2) hold, is $\mu_{\mathrm{c}}$ an integral varifold?}
Under (Q1), the tangent plane of $\|V\|$ at $\mathcal{H}^n$-a.e.\ $p \in \partial^*\Omega(x_0)$ coincides with $T_p\partial^*\Omega(x_0)$, and the sheet decomposition (Lemma~\ref{lem:sheet-decomposition}) applies there. Under (Q1)--(Q2), the decomposition also applies at $\mathcal{H}^n$-a.e.\ point of $\spt \mu_{\mathrm{c}}$, giving integer multiplicity by the graphical sheet count.

Under the multiplicity-one conjecture, mass cancellation does not occur; then $\mu_{\mathrm{c}} = 0$, Lemma~\ref{lem:density-lower-bound} gives density $\geq 1$ on all of $\spt\|V\| = \partial^*\Omega(x_0)$, and integrality follows. The gap in (Q2) and the multiplicity-one question thus form a parallel pair of open problems in the mass cancellation case.
\end{remark}

\begin{lemma}[Sheet decomposition]\label{lem:sheet-decomposition}
Let $(M^{n+1}, g)$ be a closed Riemannian manifold with $n \geq 2$. Let $\{\Omega_i\}$ be a nested sequence of Caccioppoli sets ($\Omega_1 \subset \Omega_2 \subset \cdots$) with $|\partial^*\Omega_i| \to V$ as varifolds, where $V$ is a rectifiable stationary $n$-varifold with support $\Sigma$. Let $p \in \Sigma$ be a point with unique tangent plane $P = T_p\Sigma$ and positive density $\theta(p) = \Theta(\|V\|, p) > 0$ (by rectifiability, these conditions hold at $\mathcal{H}^n$-a.e.\ point of $\Sigma$; see Section~\ref{sec:p2-prelim}).

Then for every $\varepsilon > 0$, there exist $r_0 > 0$ and $i_0 \in \mathbb{N}$ such that for a.e.\ $r \in (0, r_0)$ and all $i \geq i_0$, the following holds:
\begin{enumerate}[label=\textup{(\roman*)}]
    \item $\partial^*\Omega_i \cap B_r(p)$ decomposes into exactly $k$ connected graphical components $\Sigma_i^1, \ldots, \Sigma_i^k$, where $k = k(r)$ is independent of $i$ and satisfies $\theta(p) = k$;
    \item each $\Sigma_i^j$ is the graph of a Lipschitz function $u_i^j \colon \Omega_i^j \to \mathbb{R}$ over a domain $\Omega_i^j \subset D_r$ with Lipschitz constant $\mathrm{Lip}(u_i^j) < \varepsilon$ and $\mathcal{H}^n(D_r \setminus \Omega_i^j) < \varepsilon \omega_n r^n$, where $D_r = P \cap B_r(p)$;
    \item the area of each sheet satisfies $|\mathcal{H}^n(\Sigma_i^j) - \omega_n r^n| < \varepsilon \omega_n r^n$.
\end{enumerate}
In particular, $\theta(p) \in \mathbb{Z}_{>0}$, so $V$ is an integral $n$-varifold (Definition~\ref{def:integral-varifold}).
\end{lemma}

\begin{figure}[ht]
    \centering
    \begin{minipage}[b]{0.58\textwidth}
      \centering
      \begin{tikzpicture}[xshift=-0.8cm]
        % Junction point
        \coordinate (J) at (0,0);
        % Long arm: slight upward curve to the right
        \draw[thick] (J) to[out=10, in=170] coordinate[pos=0.75] (P) (3, 0.4);
        % Tangent line at P (blue)
        \draw[thick, bindB] ($(P) + (-1.2, -0.12)$) node[left] {$P = T_p\Sigma$} -- ($(P) + (1.2, 0.12)$);
        % Two short arms going upper-left and lower-left
        \draw[thick] (J) to[out=140, in=0] (-1.8, 1.2);
        \draw[thick] (J) to[out=220, in=0] (-1.8, -1.2);
        % Junction point dot
        \fill (J) circle (2pt);
        % Red sheets — thin bright red (outside circle)
        \draw[thin, red] (-2.2, 1.6) to[out=-10, in=170] ($(P) + (-0.8, 0.3)$) to[out=-10, in=170] ($(P) + (0.8, 0.25)$) to[out=-10, in=170] (3.5, 0.6);
        \draw[thin, red] (-2.2, -1.6) to[out=10, in=190] ($(P) + (-0.8, -0.25)$) to[out=-5, in=175] ($(P) + (0.8, -0.25)$) to[out=-5, in=175] (3.5, -0.3);
        \draw[thin, red] (-1.5, 1.0) to[out=350, in=90] (-0.8, 0) to[out=270, in=10] (-1.5, -1.0);
        % Red sheets — thick (inside circle, clipped)
        \begin{scope}
          \clip (P) circle (0.8);
          \draw[very thick, bindA] (-2.2, 1.6) to[out=-10, in=170] ($(P) + (-0.8, 0.3)$) to[out=-10, in=170] coordinate[pos=0.5] (S1) ($(P) + (0.8, 0.25)$) to[out=-10, in=170] (3.5, 0.6);
          \draw[very thick, bindA] (-2.2, -1.6) to[out=10, in=190] ($(P) + (-0.8, -0.25)$) to[out=-5, in=175] coordinate[pos=0.5] (S2) ($(P) + (0.8, -0.25)$) to[out=-5, in=175] (3.5, -0.3);
        \end{scope}
        % Cut curves at B_r(p) boundary: white eraser then gray circle
        \draw[white, line width=4pt] (P) circle (0.8);
        \draw[gray, thick] (P) circle (0.8);
        \node[gray, above right] at ($(P) + (0.55, 0.55)$) {$B_r(p)$};
        % Blue point on the long arm (drawn on top)
        \fill[bindB] (P) circle (2.5pt);
        % Labels for the two connected components — on top layer
        \node[bindA] (L1) at ($(P) + (-0.2, 1.4)$) {$\Sigma_i^1$};
        \draw[bindA, ->] (L1) -- (S1);
        \node[bindA] (L2) at ($(P) + (-0.2, -1.4)$) {$\Sigma_i^2$};
        \draw[bindA, ->] (L2) -- (S2);
        % Label for all three red curves
        \node[red] at (-2.8, 0) {$\partial^*\Omega(x_i)$};
      \end{tikzpicture}\\[4pt]
      \textbf{(a)}
    \end{minipage}%
    \hfill
    \begin{minipage}[b]{0.38\textwidth}
      \centering
      \begin{tikzpicture}[scale=1.5]
        % Define P at origin for this zoomed-in view
        \coordinate (P) at (0,0);
        % Named paths for intersection computation
        \path[name path=circ] (P) circle (0.8);
        \path[name path=sheet1] (-1.5, 0.5) to[out=-3, in=183] (1.5, 0.44);
        \path[name path=sheet2] (-1.5, -0.62) to[out=3, in=177] (1.5, -0.68);
        % Find intersections
        \path[name intersections={of=circ and sheet1, by={I1L,I1R}}];
        \path[name intersections={of=circ and sheet2, by={I2L,I2R}}];
        % Gray circle
        \draw[gray, thick] (P) circle (0.8);
        \node[gray, above right] at ($(P) + (0.55, 0.55)$) {$B_r(p)$};
        % Blue tangent line (clipped to circle) — horizontal
        \begin{scope}
          \clip (P) circle (0.8);
          \draw[thick, bindB] ($(P) + (-1.2, 0)$) -- ($(P) + (1.2, 0)$);
        \end{scope}
        % Two thick red sheets (clipped to circle)
        \begin{scope}
          \clip (P) circle (0.8);
          \draw[very thick, bindA] (-1.5, 0.5) to[out=-3, in=183] coordinate[pos=0.5] (S1b) (1.5, 0.44);
          \draw[very thick, bindA] (-1.5, -0.62) to[out=3, in=177] coordinate[pos=0.5] (S2b) (1.5, -0.68);
        \end{scope}
        % Vertical dashed red lines from exact intersections to D_r
        \draw[thin, bindA, dashed] (I1L) -- (I1L |- P);
        \draw[thin, bindA, dashed] (I1R) -- (I1R |- P);
        \draw[thin, bindA, dashed] (I2L) -- (I2L |- P);
        \draw[thin, bindA, dashed] (I2R) -- (I2R |- P);
        % Braces for domains on D_r
        \draw[red, thick, decorate, decoration={brace, mirror, amplitude=4pt}]
          (I1L |- P) -- (I1R |- P);
        \node[red] (OL2) at (1.2, -0.25) {$\Omega_i^2$};
        \draw[red, thin, ->] (OL2) to[bend left=40] (0.15, -0.07);
        \draw[red, thick, decorate, decoration={brace, amplitude=4pt}]
          (I2R |- P) -- (I2L |- P);
        \node[red] (OL1) at (-1.2, 0.25) {$\Omega_i^1$};
        \draw[red, thin, ->] (OL1) to[bend left=40] (-0.15, 0.07);
        % Labels outside circle with arrows
        \node[bindA] (L1b) at ($(P) + (-0.2, 1.3)$) {$\mathrm{graph}(u_i^1)$};
        \draw[bindA, ->] (L1b) -- (S1b);
        \node[bindA] (L2b) at ($(P) + (-0.2, -1.3)$) {$\mathrm{graph}(u_i^2)$};
        \draw[bindA, ->] (L2b) -- (S2b);
        % Tangent plane label
        \node[bindB, right] at ($(P) + (0.85, 0)$) {$D_r$};
      \end{tikzpicture}\\[4pt]
      \textbf{(b)}
    \end{minipage}
    \caption{Sheet decomposition (Lemma~\ref{lem:sheet-decomposition}) with $\theta(p) = k = 2$. \textbf{(a)}~Near a regular point $p$, the ball $B_r(p)$ cuts the reduced boundary $\partial^*\Omega(x_i)$ into $k = 2$ connected components $\Sigma_i^1, \Sigma_i^2$. \textbf{(b)}~Inside $B_r(p)$: each $\Sigma_i^j = \mathrm{graph}(u_i^j)$ over a domain $\Omega_i^j \subset D_r = P \cap B_r(p)$.}
    \label{fig:sheet-decomposition}
\end{figure}

Before giving the proof, we discuss the geometric content of the three conditions; see Figure~\ref{fig:sheet-decomposition} for an illustration with $k = 2$.

Condition~(i) says that, at sufficiently small scales, $\partial^*\Omega_i \cap B_r(p)$ separates into exactly $k = \theta(p)$ graphical sheets (Figure~\ref{fig:sheet-decomposition}(a)). This is the key step that converts the integer density $\theta(p)$ into a geometric count of sheets.

Condition~(ii) encodes two further properties. The Lipschitz bound $\mathrm{Lip}(u_i^j) < \varepsilon$ controls the \emph{tilt}: each sheet is nearly parallel to the approximate tangent plane $P = T_p\Sigma$, which is the graphical manifestation of small tilt-excess (Definition~\ref{def:p2-tilt-excess}). The coverage bound $\mathcal{H}^n(D_r \setminus \Omega_i^j) < \varepsilon \omega_n r^n$ says that the domain $\Omega_i^j \subset D_r$ covers all but an $\varepsilon$-fraction of the disk $D_r = P \cap B_r(p)$ (Figure~\ref{fig:sheet-decomposition}(b)). Together, the two bounds say that each sheet is a nearly complete, nearly horizontal graph over $D_r$. Figure~\ref{fig:condition-ii} illustrates two configurations that condition~(ii) excludes: a sheet with large tilt (Figure~\ref{fig:condition-ii}(a); which cannot persist as $i \to \infty$, since varifold convergence forces the tilt-excess to decay), and a sheet with poor coverage (Figure~\ref{fig:condition-ii}(b); whose small domain is incompatible with the mass convergence $\|\partial^*\Omega_i\|(B_r(p)) \to \theta(p)\omega_n r^n$).

\begin{figure}[ht]
    \centering
    \begin{minipage}[b]{0.45\textwidth}
      \centering
      \begin{tikzpicture}[scale=1.3]
        \coordinate (O) at (0,0);
        % Named paths for intersection
        \path[name path=circA] (O) circle (0.8);
        \path[name path=sheetA] (-1.5, -0.15) -- (1.5, 0.75);
        % Gray circle
        \draw[gray, thick] (O) circle (0.8);
        % D_r (blue, horizontal)
        \begin{scope}
          \clip (O) circle (0.8);
          \draw[thick, bindB] (-1.2, 0) -- (1.2, 0);
        \end{scope}
        \node[bindB, right] at (0.85, 0) {$D_r$};
        % Tilted sheet above D_r (large tilt)
        \begin{scope}
          \clip (O) circle (0.8);
          \draw[very thick, bindA] (-1.5, -0.15) -- (1.5, 0.75);
        \end{scope}
        \node[bindA, above right] at (0.3, 0.75) {$\Sigma_i^j$};
        % Intersection points
        \path[name intersections={of=circA and sheetA, by={IAL,IAR}}];
        % Dashed vertical projections to D_r
        \draw[thin, bindA, dashed] (IAL) -- (IAL |- O);
        \draw[thin, bindA, dashed] (IAR) -- (IAR |- O);
        % Coverage domain brace
        \draw[red, thick, decorate, decoration={brace, amplitude=3pt}]
          (IAL |- {(0,-0.1)}) -- (IAR |- {(0,-0.1)}) node[midway, below=4pt, red] {$\Omega_i^j$};
      \end{tikzpicture}\\[4pt]
      \textbf{(a)} Large tilt
    \end{minipage}%
    \hfill
    \begin{minipage}[b]{0.45\textwidth}
      \centering
      \begin{tikzpicture}[scale=1.3]
        \coordinate (O) at (0,0);
        % Named paths for intersection
        \path[name path=circ] (O) circle (0.8);
        \path[name path=sheet] (-1.2, 0.75) -- (1.2, 0.75);
        % Gray circle
        \draw[gray, thick] (O) circle (0.8);
        % D_r (blue, horizontal)
        \begin{scope}
          \clip (O) circle (0.8);
          \draw[thick, bindB] (-1.2, 0) -- (1.2, 0);
        \end{scope}
        \node[bindB, right] at (0.85, 0) {$D_r$};
        % Short horizontal sheet high above D_r (poor coverage)
        \begin{scope}
          \clip (O) circle (0.8);
          \draw[very thick, bindA] (-1.2, 0.75) -- (1.2, 0.75);
        \end{scope}
        \node[bindA, above] at (0, 0.85) {$\Sigma_i^j$};
        % Intersection points
        \path[name intersections={of=circ and sheet, by={IL,IR}}];
        % Dashed vertical projections to D_r
        \draw[thin, bindA, dashed] (IL) -- (IL |- O);
        \draw[thin, bindA, dashed] (IR) -- (IR |- O);
        % Coverage domain brace
        \draw[red, thick, decorate, decoration={brace, amplitude=3pt}]
          (IL |- {(0,-0.1)}) -- (IR |- {(0,-0.1)}) node[midway, below=4pt, red] {$\Omega_i^j$};
      \end{tikzpicture}\\[4pt]
      \textbf{(b)} Poor coverage
    \end{minipage}
    \caption{Two configurations excluded by condition~(ii) of Lemma~\ref{lem:sheet-decomposition}. In each case the domain $\Omega_i^j$ (brace) is the orthogonal projection of $\Sigma_i^j$ onto $D_r$.}
    \label{fig:condition-ii}
\end{figure}

Condition~(iii) is not independent: it follows from conditions~(i) and~(ii) by the area formula for graphs. Since each sheet $\Sigma_i^j = \mathrm{graph}(u_i^j)$, we have
\[
    \mathcal{H}^n(\Sigma_i^j) = \int_{\Omega_i^j} \sqrt{1 + |Du_i^j|^2} \, d\mathcal{H}^n.
\]
The Lipschitz bound $\mathrm{Lip}(u_i^j) < \varepsilon$ from condition~(ii) gives $\sqrt{1 + |Du_i^j|^2} \in (1, \sqrt{1+\varepsilon^2})$, so
\[
    \mathcal{H}^n(\Omega_i^j) \leq \mathcal{H}^n(\Sigma_i^j) \leq \sqrt{1+\varepsilon^2} \cdot \mathcal{H}^n(\Omega_i^j).
\]
The coverage bound $\mathcal{H}^n(D_r \setminus \Omega_i^j) < \varepsilon\omega_n r^n$ gives $\mathcal{H}^n(\Omega_i^j) \geq (1-\varepsilon)\omega_n r^n$. Combining:
\[
    (1-\varepsilon)\omega_n r^n \leq \mathcal{H}^n(\Sigma_i^j) \leq \sqrt{1+\varepsilon^2} \cdot \omega_n r^n.
\]
For small $\varepsilon$, both sides are close to $\omega_n r^n$, yielding $|\mathcal{H}^n(\Sigma_i^j) - \omega_n r^n| < C\varepsilon \cdot \omega_n r^n$. We include condition~(iii) in the lemma statement for convenience of citation, but it is logically redundant.

We now turn to the proof of Lemma~\ref{lem:sheet-decomposition}.

\begin{proof}
Working in normal coordinates centered at $p$, we identify a neighborhood of $p$ with a neighborhood of the origin in $\mathbb{R}^{n+1}$, the tangent plane $P$ with $\mathbb{R}^n \times \{0\}$, and write $D_r = \{y \in \mathbb{R}^n : |y| < r\}$ for the $n$-disk in $P$. By definition of density (Definition~\ref{def:p2-density}),
\begin{equation}\label{eq:density-ratio-limit}
    \omega_n^{-1} r^{-n} \|V\|(B_r(p)) \to \Theta(\|V\|, p) = \theta(p) \quad \text{as } r \to 0,
\end{equation}
We claim that for any $\varepsilon > 0$, all $r > 0$ sufficiently small, and all $i$ sufficiently large (depending on $r$), the surface $\partial^*\Omega_i \cap B_{2r}(p)$ has mass at most $C\omega_n r^n$ and tilt-excess $E(|\partial^*\Omega_i|, B_{2r}(p), P) < \varepsilon$. The mass bound follows from varifold convergence $|\partial^*\Omega_i| \overset{*}{\rightharpoonup} V$ and~\eqref{eq:density-ratio-limit}. For the tilt-excess: Lemma~\ref{lem:excess-vanishing} gives $E(V, B_{2r}(p), P) \to 0$ as $r \to 0$. Since varifold convergence is weak-$*$ convergence of Radon measures on the Grassmann bundle $M \times G(n+1,n)$, the continuous integrand $|S - P|^2$ yields
\begin{equation}\label{eq:tilt-excess-transfer}
    E(|\partial^*\Omega_i|, B_{2r}(p), P) \to E(V, B_{2r}(p), P) \quad \text{as } i \to \infty
\end{equation}
for a.e.\ $r$ (those satisfying $\|V\|(\partial B_{2r}(p)) = 0$). Choosing first $r$ small, then $i$ large, gives the claim. The surface $\partial^*\Omega_i \cap B_{2r}(p)$ therefore decomposes into finitely many Lipschitz graphs $\Sigma_i^1, \ldots, \Sigma_i^{k_i}$ over $P$ with $\mathrm{Lip}(u_i^j) < C\varepsilon^{1/4}$, plus a non-graphical part of $\mathcal{H}^n$-measure $\leq C\sqrt{\varepsilon}\,r^n$.

It remains to show that each sheet nearly fills $D_r$ (condition~(ii)) and that the sheet count stabilizes (condition~(i)). The strategy is as follows. By the sphere constraint (Step~1), each sheet has height bounded by $r$ and domain containing $D_{r/2}$, giving a crude area lower bound. By the tangent plane condition (Step~2), this crude bound is upgraded to a sharp height estimate $|u_i^j| \leq C\varepsilon^{1/4} r$, which in turn gives sharp coverage $\Omega_i^j \supset D_{r(1-C\varepsilon^{1/2})}$. The sharp coverage is needed for the counting argument: each sheet has area $(1 \pm C\varepsilon)\omega_n r^n$, so $k_i = \theta(p)$.
\medskip\noindent\textbf{Step~1} (Sphere constraint and crude bound). Since $\partial^*\Omega_i$ has no boundary, each graph $\Sigma_i^j$ can only terminate at $\partial B_r(p)$, so every boundary point of the domain $\Omega_i^j$ satisfies $|y|^2 + |u_i^j(y)|^2 = r^2$. Consider a sheet that intersects $B_{r/2}(p)$: there exists a point with $|u_i^j(y_0)| < r/2$, and since $\mathrm{Lip}(u_i^j) < C\varepsilon^{1/4}$, we have $|u_i^j| < r/2 + C\varepsilon^{1/4} r$ on $\Omega_i^j$. Substituting into the sphere constraint, $\Omega_i^j \supset D_{r/2}$ and hence $\mathcal{H}^n(\Omega_i^j) \geq 2^{-n}\omega_n r^n$.

\medskip\noindent\textbf{Step~2} (Tangent plane upgrade). The goal is to show $|u_i^j| \leq C\varepsilon^{1/4} r$ everywhere on each contributing sheet, upgrading the crude bound $|u_i^j| < r$ from Step~1. The key input is that the tangent varifold of $V$ at $p$ is $\theta(p)|P|$, which is supported on $P$. This means that the mass of $V$ far from $P$ is negligible at scale $r$:
\[
    \|V\|(B_r(p) \cap \{\mathrm{dist}(\cdot, P) > \varepsilon^{1/4} r\}) = o(r^n) \quad \text{as } r \to 0.
\]
By varifold convergence, the same holds for $\partial^*\Omega_i$ (for $r$ small, then $i$ large). Since each contributing sheet has area $\geq 2^{-n}\omega_n r^n$ (Step~1) and the total area of $\partial^*\Omega_i \cap B_r(p)$ beyond height $\varepsilon^{1/4} r$ from $P$ is $o(r^n)$, each sheet must contain a point with $|u_i^j| \leq \varepsilon^{1/4} r$. Propagating via $\mathrm{Lip}(u_i^j) < C\varepsilon^{1/4}$ gives $|u_i^j| \leq \varepsilon^{1/4} r + C\varepsilon^{1/4} \cdot 2r \leq C\varepsilon^{1/4} r$ everywhere. The sphere constraint then yields $\Omega_i^j \supset D_{r(1 - C\varepsilon^{1/2})}$, establishing~(ii). Condition~(iii) follows from the area formula for Lipschitz graphs.

\medskip\noindent\textbf{Step~3} (Counting). We count in $B_{r/2}(p)$. Each contributing sheet has domain $\Omega_i^j \supset D_{r/2}$ \textup{(Step~1)} and height $|u_i^j| \leq C\varepsilon^{1/4} r \ll r/2$ \textup{(Step~2)}, so the graph over $D_{r/2}$ lies inside $B_{r/2}(p)$, and the area of each sheet in $B_{r/2}(p)$ satisfies
\[
    (1 - C\varepsilon^{1/2})\omega_n (r/2)^n \leq \mathcal{H}^n(\Sigma_i^j \cap B_{r/2}(p)) \leq (1 + C\varepsilon^{1/2})\omega_n (r/2)^n.
\]
Summing over the $k_i$ sheets, we obtain
\[
    \mathcal{H}^n(\partial^*\Omega_i \cap B_{r/2}(p)) = k_i \,\omega_n (r/2)^n + O(\varepsilon^{1/2})\,(r/2)^n.
\]
Dividing by $\omega_n(r/2)^n$, we get $k_i = \|V\|(B_{r/2}(p)) / (\omega_n(r/2)^n) + O(\varepsilon^{1/2})$. Since $\|V\|(B_{r/2}(p)) / (\omega_n(r/2)^n) \to \theta(p)$ as $r \to 0$ and $k_i \in \mathbb{Z}$, for $r$ and $\varepsilon$ sufficiently small the count $k_i$ equals $\theta(p) \in \mathbb{Z}_{>0}$, establishing~(i).
\end{proof}


\begin{remark}[Role of nestedness]\label{rem:integrality-independence}
The sheet decomposition (Lemma~\ref{lem:sheet-decomposition}) uses nestedness to order the sheets and stabilize the sheet count. The graphical decomposition step (Chebyshev + Caccioppoli $\pi$-injectivity) holds for any sequence of Caccioppoli sets, but nestedness simplifies the integrality conclusion. (Our sheet decomposition is closely related to Wickramasekera's Sheeting Theorem~\cite[Theorem~3.3]{Wic14}, but applies to the \emph{approximating} Caccioppoli boundaries $\partial^*\Omega_i$ rather than the limit, avoiding the need to assume integrality or smoothness a priori.) The $\alpha$-structural verification (Theorem~\ref{thm:alpha-hypo}) uses only the Caccioppoli structure and the min-max width. The non-excessive hypothesis enters in the stability step (Proposition~\ref{prop:stability}), which invokes~\cite[Proposition~3.1]{CLS22}. Nestedness is also used in the parity analysis underlying the multiplicity-one question.
\end{remark}
